% SOURCE_AUTHORITY: sga5_fr_workpass.tex, lines 6093--6157 (LF-normalized unit).
% SOURCE_SHA256: 791F4EFFC5E02832D5D77ED03518C8156D6F07E4C8238B03545DB93D883FBB28
\subsection*{6.13}
Seja $\ell$ um número primo $\ne p$. Para todo inteiro $n\ge1$, ponhamos
\[
\Lambda_n=\mathbb Z/\ell^n.
\]
Sob as hipóteses de 6.11, seja $L=(L_n)$ um sistema projetivo de objetos de
\[
D_{\mathrm{ctf}}(G,X,\Lambda_n)
\]
tal que
\[
L_m\lten_{\Lambda_m}\Lambda_n\xrightarrow{\sim}L_n\quad (m\ge n),
\]
e seja
\[
u\in\Hom(c_1^*L,c_2^!L)
\]
um sistema projetivo de correspondências equivariantes $u_n\in\Hom(c_1^*L_n,c_2^!L_n)$ tal que, para $m\ge n$,
\[
u_m\lten_{\Lambda_m}\Lambda_n=u_n
\]
sob a identificação $L_m\lten_{\Lambda_m}\Lambda_n=L_n$ dada pelo isomorfismo anterior. Em virtude da compatibilidade dos traços locais com a extensão de escalares, os termos locais, para $g\in G$,
\[
\Tr_{\Lambda_n}(gu_n)\in H^0(X^{gc},K\Lambda_{n,X^{gc}})
\]
definem um elemento
\[
\Tr_{\mathbb Z_{\ell}}(gu)\in
H^0(X^{gc},K_{\ell,X^{gc}})
\stackrel{\mathrm{dfn}}{=}
\varprojlim_n H^0(X^{gc},K\Lambda_{n,X^{gc}}),
\]
e os elementos compatíveis
\[
f_*^{gc}\Tr_{\Lambda_n}(gu_n)=\Tr_{\Lambda_n}(gf_*u_n)
\]
definem um elemento
\[
f_*\Tr_{\mathbb Z_{\ell}}(gu)\in
H^0(Y^d,K_{\ell,Y^d})
\stackrel{\mathrm{dfn}}{=}
\varprojlim_n H^0(Y^d,K\Lambda_{n,Y^d}).
\]
Suponhamos ainda que
\[
f_*L_n\in\Ob D_{\mathrm{ctf}}({}_{\Lambda_n[G]}Y)
\]
para todo $n$. Então, os traços locais
\[
\Tr_{\Lambda_n[G]}(f_*u_n)
\]
definem igualmente um elemento
\[
\Tr_{\mathbb Z_{\ell}[G]}(f_*u)\in
H^0(Y^d,K\mathbb Z_{\ell}[G]_{Y^d,\natural})
\stackrel{\mathrm{dfn}}{=}
\varprojlim_n H^0(Y^d,K\Lambda_n[G]_{Y^d,\natural}).
\]
De 6.12 deduz-se que
\begin{equation}
f_*^{gc}\Tr_{\mathbb Z_{\ell}}(gu)
=
|Z_g|\,\Tr_{\mathbb Z_{\ell}[G]}(f_*u)(g^{-1}).
\tag{6.13.1}
\end{equation}
